
\FigureLayoutDeclare{img/2018/national_paper_1_liberal/q9_fig1.png}{9.5cm}{50bp 24bp 70bp 0bp}
\FigureLayoutDeclare{img/2018/national_paper_1_liberal/q18_fig1.png}{9.0cm}{137bp 47bp 23bp 25bp}

\chapter{2018年全国I卷（文）}

\section{单选题}

\begin{problem}
已知集合 \(A=\{0,2\}\)，\(B=\{-2,-1,0,1,2\}\)，则 \(A\cap B=\)（\quad）

\choices
    {\(\{0,2\}\)}
    {\(\{1,2\}\)}
    {\(\{0\}\)}
    {\(\{-2,-1,0,1,2\}\)}
\end{problem}

\begin{answer}
A
\end{answer}

\begin{solution}
交集中的元素必须同时属于集合 \(A\) 和集合 \(B\). 逐个比较可知，\(0\)，\(2\) 同时属于 \(A\)，\(B\)，而 \(A\) 中没有其他元素，因此 \(A\cap B=\{0,2\}\)，故选 A.
\end{solution}

\begin{problem}
设 \(z=\dfrac{1-\i}{1+\i}+2\i\)，则 \(|z|=\)（\quad）

\choices
    {\(0\)}
    {\(\dfrac12\)}
    {\(1\)}
    {\(\sqrt2\)}
\end{problem}

\begin{answer}
C
\end{answer}

\begin{solution}
分子分母同乘 \(1-\i\)，得
\[
\frac{1-\i}{1+\i}=\frac{(1-\i)^2}{(1+\i)(1-\i)}=\frac{-2\i}{2}=-\i
\]
所以 \(z=-\i+2\i=\i\)，从而 \(|z|=1\)，故选 C.
\end{solution}

\begin{problem}
某地区经过一年的新农村建设，农村的经济收入增加了一倍，实现翻番. 为更好地了解该地区农村的经济收入变化情况，统计了该地区新农村建设前后农村的经济收入构成比例，得到如下饼图：

\bitmapfigure[max width=0.92\linewidth]{img/2018/national_paper_1_liberal/q3_fig1.png}

则下面结论中不正确的是（\quad）

\choices
    {新农村建设后，种植收入减少}
    {新农村建设后，其他收入增加了一倍以上}
    {新农村建设后，养殖收入增加了一倍}
    {新农村建设后，养殖收入与第三产业收入的总和超过了经济收入的一半}
\end{problem}

\begin{answer}
A
\end{answer}

\begin{solution}
设建设前农村经济收入为 \(S\)，则建设后为 \(2S\). 建设前种植收入为 \(0.60S\)，建设后种植收入为 \(0.37\times2S=0.74S\)，所以种植收入增加而不是减少，选项 A 不正确.

建设前其他收入为 \(0.04S\)，建设后为 \(0.05\times2S=0.10S\)，后者是前者的 \(2.5\) 倍，增加了一倍以上；建设前养殖收入为 \(0.30S\)，建设后为 \(0.30\times2S=0.60S\)，恰好增加一倍；建设后养殖收入与第三产业收入的总和为 \((0.30+0.28)\times2S=1.16S>S\)，超过经济收入的一半. 因此不正确的是 A.
\end{solution}

\begin{problem}
已知椭圆 \(C:\dfrac{x^2}{a^2}+\dfrac{y^2}{4}=1\) 的一个焦点为 \((2,0)\)，则 \(C\) 的离心率为（\quad）

\choices
    {\(\dfrac13\)}
    {\(\dfrac12\)}
    {\(\dfrac{\sqrt2}{2}\)}
    {\(\dfrac{2\sqrt2}{3}\)}
\end{problem}

\begin{answer}
C
\end{answer}

\begin{solution}
椭圆的半焦距为 \(c=2\)，短半轴平方为 \(b^2=4\). 由 \(a^2=b^2+c^2\)，得 \(a^2=4+4=8\)，所以 \(a=2\sqrt2\). 因此离心率
\[
e=\frac ca=\frac{2}{2\sqrt2}=\frac{\sqrt2}{2}
\]
故选 C.
\end{solution}

\begin{problem}
已知圆柱的上、下底面的中心分别为 \(O_1\)，\(O_2\)，过直线 \(O_1O_2\) 的平面截该圆柱所得的截面是面积为 \(8\) 的正方形，则该圆柱的表面积为（\quad）

\choices
    {\(12\sqrt2\pi\)}
    {\(12\pi\)}
    {\(8\sqrt2\pi\)}
    {\(10\pi\)}
\end{problem}

\begin{answer}
B
\end{answer}

\begin{solution}
设圆柱的底面半径为 \(r\)，高为 \(h\). 过轴线的截面是矩形，其两边长分别为 \(2r\) 和 \(h\). 由题意该矩形是面积为 \(8\) 的正方形，所以正方形边长为 \(2\sqrt2\)，从而
\(2r=h=2\sqrt2\)，\(r=\sqrt2\)
圆柱的表面积为
\[
2\pi r^2+2\pi rh=2\pi\times2+2\pi\times\sqrt2\times2\sqrt2=12\pi
\]
故选 B.
\end{solution}

\begin{problem}
设函数 \(f(x)=x^3+(a-1)x^2+ax\). 若 \(f(x)\) 为奇函数，则曲线 \(y=f(x)\) 在点 \((0,0)\) 处的切线方程为（\quad）

\choices
    {\(y=-2x\)}
    {\(y=-x\)}
    {\(y=2x\)}
    {\(y=x\)}
\end{problem}

\begin{answer}
D
\end{answer}

\begin{solution}
奇函数的偶次幂项系数必须为零，因此 \(a-1=0\)，得 \(a=1\). 于是 \(f(x)=x^3+x\)，且 \(f'(x)=3x^2+1\)，所以 \(f'(0)=1\). 曲线在 \((0,0)\) 处的切线方程为 \(y=x\)，故选 D.
\end{solution}

\begin{problem}
在 \(\triangle ABC\) 中，\(AD\) 为 \(BC\) 边上的中线，\(E\) 为 \(AD\) 的中点，则 \(\overrightarrow{EB}=\)（\quad）

\choices
    {\(\dfrac34\overrightarrow{AB}-\dfrac14\overrightarrow{AC}\)}
    {\(\dfrac14\overrightarrow{AB}-\dfrac34\overrightarrow{AC}\)}
    {\(\dfrac34\overrightarrow{AB}+\dfrac14\overrightarrow{AC}\)}
    {\(\dfrac14\overrightarrow{AB}+\dfrac34\overrightarrow{AC}\)}
\end{problem}

\begin{answer}
A
\end{answer}

\begin{solution}
设 \(\overrightarrow{AB}=\bs u\)，\(\overrightarrow{AC}=\bs v\)，并取 \(A\) 为位置参考点. 由于 \(D\) 是 \(BC\) 的中点，有 \(\overrightarrow{AD}=\frac12(\bs u+\bs v)\). 又因为 \(E\) 是 \(AD\) 的中点，所以 \(\overrightarrow{AE}=\frac14(\bs u+\bs v)\).
于是
\[
\overrightarrow{EB}=\overrightarrow{AB}-\overrightarrow{AE}=\bs u-\frac14(\bs u+\bs v)=\frac34\bs u-\frac14\bs v
\]
即 \(\overrightarrow{EB}=\frac34\overrightarrow{AB}-\frac14\overrightarrow{AC}\)，故选 A.
\end{solution}

\begin{problem}
已知函数 \(f(x)=2\cos^2 x-\sin^2 x+2\)，则（\quad）

\choices
    {\(f(x)\) 的最小正周期为 \(\pi\)，最大值为 \(3\)}
    {\(f(x)\) 的最小正周期为 \(\pi\)，最大值为 \(4\)}
    {\(f(x)\) 的最小正周期为 \(2\pi\)，最大值为 \(3\)}
    {\(f(x)\) 的最小正周期为 \(2\pi\)，最大值为 \(4\)}
\end{problem}

\begin{answer}
B
\end{answer}

\begin{solution}
利用 \(\sin^2x=1-\cos^2x\)，得
\[
f(x)=2\cos^2x-(1-\cos^2x)+2=3\cos^2x+1=\frac52+\frac32\cos2x
\]
因此 \(f(x)\) 的最小正周期为 \(\pi\)，最大值为 \(\frac52+\frac32=4\)，故选 B.
\end{solution}

\begin{problem}
某圆柱的高为 \(2\)，底面周长为 \(16\)，其三视图如图. 圆柱表面上的点 \(M\) 在正视图上的对应点为 \(A\)，圆柱表面上的点 \(N\) 在左视图上的对应点为 \(B\)，则在此圆柱侧面上，从 \(M\) 到 \(N\) 的路径中，最短路径的长度为（\quad）

\bitmapfigure[max width=0.72\linewidth]{img/2018/national_paper_1_liberal/q9_fig1.png}

\choices
    {\(2\sqrt{17}\)}
    {\(2\sqrt5\)}
    {\(3\)}
    {\(2\)}
\end{problem}

\begin{answer}
B
\end{answer}

\begin{solution}
由三视图可知，点 \(M\) 与点 \(N\) 的高度分别在上、下底面所在的水平位置，所以两点的竖直距离为圆柱的高 \(2\). 同时，正视图和左视图中的对应位置表明，经过 \(M\)，\(N\) 的两条母线在底面圆周上相差 \(90^\circ\)，故沿较短圆弧的弧长为底面周长的四分之一，即 \(16/4=4\).

将圆柱侧面沿这段较短圆弧展开，得到一个矩形；矩形的高为 \(2\)，底边上两点的距离为 \(4\)，而侧面上的最短路径对应展开矩形中的线段. 因此最短路径长为 \(\sqrt{2^2+4^2}=2\sqrt5\)，故选 B.
\end{solution}

\begin{problem}
在长方体 \(ABCD-A_1B_1C_1D_1\) 中，\(AB=BC=2\)，\(AC_1\) 与平面 \(BB_1C_1C\) 所成的角为 \(30^\circ\)，则该长方体的体积为（\quad）

\choices
    {\(8\)}
    {\(6\sqrt2\)}
    {\(8\sqrt2\)}
    {\(8\sqrt3\)}
\end{problem}

\begin{answer}
C
\end{answer}

\begin{solution}
设长方体的高为 \(h\). 取 \(AC_1\) 为斜边，则 \(AC_1\) 在平面 \(BB_1C_1C\) 的法线方向上的分量为 \(AB=2\)，因为直线与平面所成角为 \(30^\circ\)，所以
\(\sin30^\circ=\frac{2}{AC_1}\)，得 \(AC_1=4\). 又 \(AC^2=AB^2+BC^2=8\)，故
\(h^2=AC_1^2-AC^2=16-8=8\)，\(h=2\sqrt2\)
所以长方体的体积为
\[
AB\cdot BC\cdot h=2\times2\times2\sqrt2=8\sqrt2
\]
故选 C.
\end{solution}

\begin{problem}
已知角 \(\alpha\) 的顶点为坐标原点，始边与 \(x\) 轴的非负半轴重合，终边上有两点 \(A(1,a)\)，\(B(2,b)\)，且 \(\cos2\alpha=\dfrac23\)，则 \(|a-b|=\)（\quad）

\choices
    {\(\dfrac15\)}
    {\(\dfrac{\sqrt5}{5}\)}
    {\(\dfrac{2\sqrt5}{5}\)}
    {\(1\)}
\end{problem}

\begin{answer}
B
\end{answer}

\begin{solution}
由于 \(A(1,a)\)，\(B(2,b)\) 在角 \(\alpha\) 的同一条终边上，且横坐标之比为 \(1:2\)，所以 \(b=2a\). 又因为点 \(A\) 的横坐标为 \(1\)，有 \(\tan\alpha=a\). 利用二倍角公式，\(\cos2\alpha=\frac{1-\tan^2\alpha}{1+\tan^2\alpha}=\frac{1-a^2}{1+a^2}=\frac23\)，解得 \(a^2=\frac15\). 于是 \(|a-b|=|a-2a|=|a|=\frac{\sqrt5}{5}\)，故选 B.
\end{solution}

\begin{problem}
设函数 \(f(x)=\begin{cases}
2^{-x},&x\le0,\\
1,&x>0
\end{cases}\)，则满足 \(f(x+1)<f(2x)\) 的 \(x\) 的取值范围是（\quad）

\choices
    {\((-\infty,-1]\)}
    {\((0,+\infty)\)}
    {\((-1,0)\)}
    {\((-\infty,0)\)}
\end{problem}

\begin{answer}
D
\end{answer}

\begin{solution}
当 \(x\leq-1\) 时，\(x+1\leq0\)，\(2x\leq0\)，所以
\(f(x+1)=2^{-x-1}\)，\(f(2x)=2^{-2x}\). 而 \(-x-1<-2x\) 等价于 \(x<1\)，故此时不等式对所有 \(x\leq-1\) 成立.
当 \(-1<x\leq0\) 时，\(x+1>0\)，\(2x\leq0\)，所以不等式化为
\(1<2^{-2x}\). 这等价于 \(-2x>0\)，即 \(-1<x<0\). 当 \(x>0\) 时，\(x+1>0\)，\(2x>0\)，两边函数值均为 \(1\)，不等式不成立. 因此 \(x\) 的取值范围为 \((-\infty,0)\)，故选 D.
\end{solution}

\section{填空题}

\begin{problem}
已知函数 \(f(x)=\log_2(x^2+a)\)，若 \(f(3)=1\)，则 \(a=\fillinblank{}\).
\end{problem}

\begin{answer}
\(-7\)
\end{answer}

\begin{solution}
由 \(f(3)=1\)，得 \(\log_2(3^2+a)=1\)，所以 \(9+a=2\)，解得 \(a=-7\). 此时 \(3^2+a=2>0\)，满足对数真数为正的条件.
\end{solution}

\begin{problem}
若 \(x\)，\(y\) 满足约束条件 \(\begin{cases}
x-2y-2\le0,\\
x-y+1\ge0,\\
y\le0
\end{cases}\)，则 \(z=3x+2y\) 的最大值为\fillinblank{}.
\end{problem}

\begin{answer}
6
\end{answer}

\begin{solution}
约束条件可写为 \(x\leq2y+2\)，\(x\geq y-1\)，\(y\leq0\).
由 \(y-1\leq2y+2\) 得 \(y\geq-3\)，所以 \(-3\leq y\leq0\). 因为目标函数 \(z=3x+2y\) 中 \(x\) 的系数为正，在固定 \(y\) 时应取 \(x=2y+2\). 于是
\[
z\leq3(2y+2)+2y=8y+6\leq6
\]
当 \(y=0\)，\(x=2\) 时取等号，所以 \(z\) 的最大值为 \(6\).
\end{solution}

\begin{problem}
直线 \(y=x+1\) 与圆 \(x^2+y^2+2y-3=0\) 交于 \(A\)，\(B\) 两点，则 \(|AB|=\fillinblank{}\).
\end{problem}

\begin{answer}
\(2\sqrt2\)
\end{answer}

\begin{solution}
将直线方程 \(y=x+1\) 代入圆方程，得 \(x^2+(x+1)^2+2(x+1)-3=0\)，即 \(2x^2+4x=0\)，所以 \(x=0\) 或 \(x=-2\). 对应的交点为 \((0,1)\)，\((-2,-1)\)，因此 \(|AB|=\sqrt{(0+2)^2+(1+1)^2}=2\sqrt2\).
\end{solution}

\begin{problem}
\(\triangle ABC\) 的内角 \(A\)，\(B\)，\(C\) 的对边分别为 \(a\)，\(b\)，\(c\). 已知 \(b\sin C+c\sin B=4a\sin B\sin C\)，\(b^2+c^2-a^2=8\)，则 \(\triangle ABC\) 的面积为\fillinblank{}.
\end{problem}

\begin{answer}
\(\dfrac{2\sqrt3}{3}\)
\end{answer}

\begin{solution}
由正弦定理 \(b=\dfrac{a\sin B}{\sin A}\)，\(c=\dfrac{a\sin C}{\sin A}\)，代入已知等式，得 \(\frac{2a\sin B\sin C}{\sin A}=4a\sin B\sin C\).
由于三角形内角的正弦均为正，所以 \(\sin A=\dfrac12\)，从而 \(A=\dfrac{\pi}{6}\) 或 \(A=\dfrac{5\pi}{6}\). 另一方面，由余弦定理
\(b^2+c^2-a^2=2bc\cos A=8\).
左式为正，故 \(\cos A>0\)，只能有 \(A=\dfrac{\pi}{6}\). 于是 \(bc=\dfrac{4}{\cos A}=\dfrac{8}{\sqrt3}\)，三角形面积为
\[
S=\frac12bc\sin A=\frac12\cdot\frac{8}{\sqrt3}\cdot\frac12=\frac{2\sqrt3}{3}
\]
\end{solution}

\section{解答题}

\begin{problem}
已知数列 \(\{a_n\}\) 满足 \(a_1=1\)，\(na_{n+1}=2(n+1)a_n\)，设 \(b_n=\dfrac{a_n}{n}\).

\begin{enumerate}
\item 求 \(b_1\)，\(b_2\)，\(b_3\)；
\item 判断数列 \(\{b_n\}\) 是否为等比数列，并说明理由；
\item 求 \(\{a_n\}\) 的通项公式.
\end{enumerate}
\end{problem}

\begin{solution}
\begin{enumerate}
\item 由 \(a_1=1\)，得 \(b_1=\dfrac{a_1}{1}=1\). 当 \(n=1\) 时，\(a_2=4\)，所以 \(b_2=\dfrac{a_2}{2}=2\). 当 \(n=2\) 时，\(a_3=3a_2=12\)，所以 \(b_3=\dfrac{a_3}{3}=4\).
\item 由已知递推式，除以 \(n(n+1)\)，得
\(\frac{a_{n+1}}{n+1}=2\frac{a_n}{n}\)，即 \(b_{n+1}=2b_n\). 因此 \(\dfrac{b_{n+1}}{b_n}=2\) 对每个正整数 \(n\) 都成立，数列 \(\{b_n\}\) 是以 \(2\) 为公比的等比数列.
\item 由 \(b_1=1\) 及公比为 \(2\)，得 \(b_n=2^{n-1}\). 又 \(a_n=nb_n\)，所以
\(a_n=n2^{n-1}\).
\end{enumerate}
\end{solution}

\begin{problem}
如图，在平行四边形 \(ABCM\) 中，\(AB=AC=3\)，\(\angle ACM=90^\circ\)，以 \(AC\) 为折痕将 \(\triangle ACM\) 折起，使点 \(M\) 到达点 \(D\) 的位置，且 \(AB\perp DA\).

\begin{enumerate}
\item 证明：平面 \(ACD\perp\) 平面 \(ABC\)；
\item \(Q\) 为线段 \(AD\) 上一点，\(P\) 为线段 \(BC\) 上一点，且 \(BP=DQ=\dfrac23DA\)，求三棱锥 \(Q-ABP\) 的体积.
\end{enumerate}

\bitmapfigure[max width=6.0cm]{img/2018/national_paper_1_liberal/q18_fig1.png}
\end{problem}

\begin{solution}
\begin{enumerate}
\item 在平行四边形 \(ABCM\) 中，\(CM\parallel AB\). 由 \(\angle ACM=90^\circ\)，得 \(AC\perp CM\)，所以 \(AC\perp AB\). 题设又给出 \(AB\perp AD\)，而 \(AC\)，\(AD\) 是平面 \(ACD\) 内相交的两条直线，因此 \(AB\perp\) 平面 \(ACD\). 由于 \(AB\) 在平面 \(ABC\) 内，故平面 \(ACD\perp\) 平面 \(ABC\).
\item 在平行四边形 \(ABCM\) 中，\(CM=AB=3\)，且 \(AC\perp CM\)，所以在直角三角形 \(ACM\) 中
\(AM=\sqrt{AC^2+CM^2}=3\sqrt2\).
折叠是绕 \(AC\) 的刚体转动，长度和与折痕的夹角保持不变，故 \(AD=AM=3\sqrt2\)，且 \(\angle DAC=\angle MAC=45^\circ\). 由 \(\angle ACD=90^\circ\)，有 \(DC\perp AC\)；又由第（1）问，\(AB\perp\) 平面 \(ACD\)，所以 \(AB\perp DC\). 因为 \(AB\)，\(AC\) 是平面 \(ABC\) 内相交的两条直线，故 \(DC\perp\) 平面 \(ABC\)，从而点 \(D\) 到平面 \(ABC\) 的距离为 \(DC=3\).

由 \(DQ=\dfrac23DA\)，得 \(AQ=\dfrac13AD\). 线段 \(AD\) 的端点 \(A\) 在平面 \(ABC\) 内，端点 \(D\) 到该平面的距离为 \(3\)，距离沿线段按比例变化，所以点 \(Q\) 到平面 \(ABC\) 的距离为 \(\dfrac13\times3=1\).

由平行四边形知 \(BC=AM=3\sqrt2\)，且
\(BP=\dfrac23DA=2\sqrt2=\dfrac23BC\). 三角形 \(ABP\) 与三角形 \(ABC\) 共有从 \(A\) 到直线 \(BC\) 的高，因此
\([ABP]=\dfrac{BP}{BC}[ABC]=\dfrac23\cdot\dfrac12\cdot AB\cdot AC=3\).
所以三棱锥 \(Q-ABP\) 的体积为
\(V=\dfrac13[ABP]\cdot1=1\).
\end{enumerate}
\end{solution}

\begin{problem}
某家庭记录了未使用节水龙头 \(50\) 天的日用水量数据（单位：\(\text{m}^3\)）和使用了节水龙头 \(50\) 天的日用水量数据，得到频数分布表如下：

未使用节水龙头 \(50\) 天的日用水量频数分布表

\begin{center}
\begin{tabular}{c|ccccccc}
\hline
日用水量 & \([0,0.1)\) & \([0.1,0.2)\) & \([0.2,0.3)\) & \([0.3,0.4)\) & \([0.4,0.5)\) & \([0.5,0.6)\) & \([0.6,0.7)\) \\
\hline
频数 & \(1\) & \(3\) & \(2\) & \(4\) & \(9\) & \(26\) & \(5\) \\
\hline
\end{tabular}
\end{center}

使用了节水龙头 \(50\) 天的日用水量频数分布表

\begin{center}
\begin{tabular}{c|cccccc}
\hline
日用水量 & \([0,0.1)\) & \([0.1,0.2)\) & \([0.2,0.3)\) & \([0.3,0.4)\) & \([0.4,0.5)\) & \([0.5,0.6)\) \\
\hline
频数 & \(1\) & \(5\) & \(13\) & \(10\) & \(16\) & \(5\) \\
\hline
\end{tabular}
\end{center}

\begin{enumerate}
\item 作出使用了节水龙头 \(50\) 天的日用水量数据的频率分布直方图；

\item 估计该家庭使用节水龙头后，日用水量小于 \(0.35\text{ m}^3\) 的概率；
\item 估计该家庭使用节水龙头后，一年能节省多少水？（一年按 \(365\) 天计算，同一组中的数据以这组数据所在区间中点的值作代表）
\end{enumerate}

\problemresetlayout
\bitmapfigure[float=inline,width=0.50\linewidth]{img/2018/national_paper_1_liberal/q19_fig1.png}
\FloatBarrier
\end{problem}

\begin{solution}
\begin{enumerate}
\item 每组组距都是 \(0.1\)，频率分布直方图的纵坐标为频率除以组距. 使用节水龙头后的第一组频率密度为 \(\frac{1/50}{0.1}=0.2\)；第二组至第六组的频率密度分别为 \(\dfrac{5/50}{0.1}=1\)，\(\dfrac{13/50}{0.1}=2.6\)，\(\dfrac{10/50}{0.1}=2\)，\(\dfrac{16/50}{0.1}=3.2\)，\(\dfrac{5/50}{0.1}=1\).
所以直方图中六个小矩形的底分别为 \([0,0.1)\)，\([0.1,0.2)\)，\([0.2,0.3)\)，\([0.3,0.4)\)，\([0.4,0.5)\)，\([0.5,0.6)\)，高依次为 \(0.2\)，\(1\)，\(2.6\)，\(2\)，\(3.2\)，\(1\).
\item 日用水量小于 \(0.35\mathrm{~m}^3\) 的天数可估计为前3组的全部天数加上第4组的一半，即
\(1+5+13+\frac12\times10=24\)，
所以所求概率约为 \(\dfrac{24}{50}=0.48\).
\item 未使用节水龙头时，以各组区间中点作为代表，每日平均用水量约为
\[
\frac{0.05\times1+0.15\times3+0.25\times2+0.35\times4+0.45\times9+0.55\times26+0.65\times5}{50}=0.48\mathrm{~m}^3
\]
使用节水龙头后，每日平均用水量约为
\[
\frac{0.05\times1+0.15\times5+0.25\times13+0.35\times10+0.45\times16+0.55\times5}{50}=0.35\mathrm{~m}^3
\]
所以每天约节省 \(0.48-0.35=0.13\mathrm{~m}^3\)，一年约节省 \(365\times0.13=47.45\mathrm{~m}^3\).
\end{enumerate}
\end{solution}

\begin{problem}
设抛物线 \(C:y^2=2x\)，点 \(A(2,0)\)，\(B(-2,0)\)，过点 \(A\) 的直线 \(l\) 与 \(C\) 交于 \(M\)，\(N\) 两点.

\begin{enumerate}
\item 当 \(l\) 与 \(x\) 轴垂直时，求直线 \(BM\) 的方程；
\item 证明：\(\angle ABM=\angle ABN\).
\end{enumerate}
\end{problem}

\begin{solution}
\begin{enumerate}
\item 当 \(l\) 与 \(x\) 轴垂直时，\(l\) 的方程为 \(x=2\). 代入抛物线方程得 \(y^2=4\)，所以 \(M\)，\(N\) 为 \((2,2)\)，\((2,-2)\) 两点. 若 \(M=(2,2)\)，则
\(k_{BM}=\frac{2-0}{2-(-2)}=\frac12\)，从而直线 \(BM\) 的方程为 \(y=\dfrac12(x+2)\)，即 \(y=\dfrac12x+1\). 若 \(M=(2,-2)\)，则 \(k_{BM}=\frac{-2-0}{2-(-2)}=-\frac12\)，
从而直线 \(BM\) 的方程为 \(y=-\dfrac12(x+2)\)，即 \(y=-\dfrac12x-1\).
\item 当 \(l\) 不与 \(x\) 轴垂直时，设其斜率为 \(k\)，则 \(l\) 的方程为 \(y=k(x-2)\). 由题设直线与抛物线交于两个不同点，故 \(k\neq0\). 令 \(M(x_1,y_1)\)，\(N(x_2,y_2)\)，则
\(x_i=\frac{y_i}{k}+2\)，
代入 \(y^2=2x\)，得 \(ky_i^2-2y_i-4k=0\). 所以
\(y_1+y_2=\frac2k\)，\(y_1y_2=-4\).
因为 \(x_i=\dfrac{y_i^2}{2}\geq0\)，所以 \(x_i+2>0\). 直线 \(BM\)，\(BN\) 的斜率分别为
\(k_{BM}=\frac{y_1}{x_1+2}=\frac{ky_1}{y_1+4k}\)，\(k_{BN}=\frac{y_2}{x_2+2}=\frac{ky_2}{y_2+4k}\).
于是
\[
k_{BM}+k_{BN}
=\frac{k\bigl(2y_1y_2+4k(y_1+y_2)\bigr)}
{y_1y_2+4k(y_1+y_2)+16k^2}=0
\]
分母不为零是因为 \(x_i+2>0\). 又由 \(y_1y_2=-4<0\) 知 \(y_1\)，\(y_2\) 异号，不妨设 \(y_1>0>y_2\)，于是 \(BM\)，\(BN\) 分别位于 \(BA\) 的两侧，且 \(k_{BM}>0>k_{BN}\). 由于 \(BA\) 是 \(x\) 轴正向，有 \(\tan\angle ABM=k_{BM}\)，\(\tan\angle ABN=-k_{BN}\)，再结合 \(k_{BM}+k_{BN}=0\)，可得 \(\angle ABM=\angle ABN\). 当 \(l\) 与 \(x\) 轴垂直时，\(M\)，\(N\) 关于 \(BA\) 对称，结论同样成立.
\end{enumerate}
\end{solution}

\begin{problem}
已知函数 \(f(x)=a\e^x-\ln x-1\).

\begin{enumerate}
\item 设 \(x=2\) 是 \(f(x)\) 的极值点，求 \(a\)，并求 \(f(x)\) 的单调区间；
\item 证明：当 \(a\ge\dfrac1\e\) 时，\(f(x)\ge0\).
\end{enumerate}
\end{problem}

\begin{solution}
\begin{enumerate}
\item 函数定义域为 \(x>0\)，且
\(f'(x)=a\e^x-\frac1x\).
因为 \(x=2\) 是极值点，所以 \(f'(2)=0\)，得
\(a\e^2-\frac12=0\)，\(a=\frac{1}{2\e^2}\)
此时
\[
f'(x)=\frac{\e^x}{2\e^2}-\frac1x
=\frac{x\e^x-2\e^2}{2\e^2x}
\]
令 \(g(x)=x\e^x\)，则 \(g'(x)=\e^x(x+1)>0\)，所以 \(g(x)\) 在 \((0,+\infty)\) 上单调递增，且 \(g(2)=2\e^2\). 因此 \(f'(x)<0\) 时 \(0<x<2\)，\(f'(x)>0\) 时 \(x>2\). 故 \(f(x)\) 的单调减区间为 \((0,2)\)，单调增区间为 \((2,+\infty)\).
\item 当 \(a\geq\dfrac1\e\) 时，对任意 \(x>0\)，有
\(f(x)=a\e^x-\ln x-1\geq\e^{x-1}-\ln x-1\).
由 \( \e^u\geq1+u\)（\(u\in\R\)）和 \(\ln x\leq x-1\)（\(x>0\)），分别取 \(u=x-1\)，可得
\[
\e^{x-1}\geq x\geq1+\ln x
\]
所以 \(\e^{x-1}-\ln x-1\geq0\)，从而 \(f(x)\geq0\).
\end{enumerate}
\end{solution}

\section{选考题：解答题}

\begin{problem}
在直角坐标系 \(xOy\) 中，曲线 \(C_1\) 的方程为 \(y=k|x|+2\). 以坐标原点为极点，\(x\) 轴正半轴为极轴建立极坐标系，曲线 \(C_2\) 的极坐标方程为 \(\rho^2+2\rho\cos\theta-3=0\).

\begin{enumerate}
\item 求 \(C_2\) 的直角坐标方程；
\item 若 \(C_1\) 与 \(C_2\) 有且仅有三个公共点，求 \(C_1\) 的方程.
\end{enumerate}
\end{problem}

\begin{solution}
\begin{enumerate}
\item 由极坐标与直角坐标的关系 \(\rho^2=x^2+y^2\)，\(\rho\cos\theta=x\)，得 \(x^2+y^2+2x-3=0\)，
所以 \(C_2\) 的直角坐标方程为 \(x^2+y^2+2x-3=0\).
\item 当 \(x=0\) 时，\(C_1\) 上的点为 \((0,2)\)，但 \(0^2+2^2+2\cdot0-3=1\neq0\)，故公共点不在 \(y\) 轴上. 令 \(t=|x|>0\). 当 \(x\geq0\) 时，\(y=kt+2\)，代入 \(C_2\) 得
\[
(1+k^2)t^2+(4k+2)t+1=0
\]
当 \(x<0\) 时，仍有 \(y=kt+2\)，代入 \(C_2\) 得
\[
(1+k^2)t^2+(4k-2)t+1=0
\]
两式的常数项和二次项系数均为正. 第一式有两个正根的条件是 \(4k+2<0\) 且判别式 \(\Delta_+=4k(3k+4)>0\) 成立，即 \(k<-\dfrac43\)；当 \(k=-\dfrac43\) 时，第一式有一个正的二重根；其余情况下第一式没有正根.

第二式有两个正根的条件是 \(4k-2<0\) 且判别式 \(\Delta_-=4k(3k-4)>0\)，这两个条件合起来等价于 \(k<0\)；当 \(k=0\) 时，第二式为 \((t-1)^2=0\)，只有一个正的二重根；当 \(k>0\) 时没有正根. 二重根只对应一个公共点，因此公共点总数为 \(3\) 的唯一情形是第一式在 \(k=-\dfrac43\) 时有一个正的二重根、第二式有两个正根. 因此 \(C_1\) 的方程为 \(y=-\dfrac43|x|+2\).
\end{enumerate}
\end{solution}

\begin{problem}
已知 \(f(x)=|x+1|-|ax-1|\).

\begin{enumerate}
\item 当 \(a=1\) 时，求不等式 \(f(x)>1\) 的解集；
\item 若 \(x\in(0,1)\) 时不等式 \(f(x)>x\) 成立，求 \(a\) 的取值范围.
\end{enumerate}
\end{problem}

\begin{solution}
\begin{enumerate}
\item 当 \(a=1\) 时，原不等式为
\(|x+1|-|x-1|>1\).
当 \(x<-1\) 时，左边等于 \(-2\)，不满足不等式；当 \(-1\leq x<1\) 时，左边等于 \(2x\)，故需 \(x>\dfrac12\)；当 \(x\geq1\) 时，左边等于 \(2\)，恒大于 \(1\). 合并得解集为
\(\left(\frac12,+\infty\right)\).
\item 因为 \(x\in(0,1)\)，所以 \(x+1>0\)，原不等式等价于
\(x+1-|ax-1|>x\)，即 \(|ax-1|<1\). 这又等价于 \(-1<ax-1<1\)，即 \(0<ax<2\).
要使它对一切 \(x\in(0,1)\) 成立，必须有 \(a>0\). 当 \(0<a\leq2\) 时，\(0<ax<a\leq2\)，且因 \(x<1\)，即使 \(a=2\) 也有 \(ax<2\)，所以不等式恒成立；若 \(a>2\)，取 \(x\in(2/a,1)\)，则 \(ax>2\)，不等式不成立. 因此 \(a\) 的取值范围为 \(0<a\leq2\).
\end{enumerate}
\end{solution}
